La sonda \textit{Juno} descriu una òrbita polar al voltant del planeta Júpiter des del dia 5 de juliol de 2016. La seva missió és estudiar l'atmosfera, l'origen i l'estructura de Júpiter, així com la seva evolució dins del Sistema Solar. Suposeu que l'òrbita és circular i que l'altura de l'òrbita sobre el planeta és de 4\,300~km.

\figura[0.4]{fig1.png}

\begin{apartat}{1}
Calculeu l'energia cinètica de \textit{Juno} i el seu període de rotació.
\end{apartat}

\begin{apartat}{1}
Calculeu l'energia que caldria comunicar-li perquè abandonés el camp gravitatori de Júpiter.
\end{apartat}

\dades{%
  Massa de Júpiter, $M_\mathrm{J} = 1,90\times 10^{27}\ \mathrm{kg}$.\\
  Radi de Júpiter, $R_\mathrm{J} = 69\,911\ \mathrm{km}$.\\
  Massa de la sonda \textit{Juno}, $m_\mathrm{Juno} = 3\,625\ \mathrm{kg}$.\\
  $G = 6,67\times 10^{-11}\ \mathrm{N\,m^2\,kg^{-2}}$.}
